% Generated draft: requires mathematical and source review.
\documentclass[11pt]{article}
\usepackage[margin=25mm]{geometry}
\usepackage{fontspec}
\setmainfont{Latin Modern Roman}
\usepackage{amsmath,amssymb,mathtools,unicode-math}
\setmathfont{Latin Modern Math}
\usepackage{enumitem,needspace,xurl,hyperref}
\setlength{\parindent}{0pt}
\setlength{\parskip}{5pt}
\setlength{\emergencystretch}{3em}
\newcommand{\N}{\mathbb{N}}
\newcommand{\Z}{\mathbb{Z}}
\newcommand{\Q}{\mathbb{Q}}
\newcommand{\R}{\mathbb{R}}
\newcommand{\C}{\mathbb{C}}
\newcommand{\F}{\mathbb{F}}
\newcommand{\A}{\mathbb{A}}
\newcommand{\PP}{\mathbb{P}}
\newcommand{\E}{\mathbb{E}}
\newcommand{\eps}{\varepsilon}
\newcommand{\dd}{\,\mathrm{d}}
\begin{document}
\section*{UnsolvedMath ID: 30002343}
\textbf{Source problem code:} catalogue entry
\subsection{Definitions and mathematical statement}
The supplied record concerns stable log surfaces $(X,\Delta)$ and the Gorenstein locus
\[
U=\{x\in X:\mathcal O_{X,x}\text{ is Gorenstein}\}.
\]
In the usual terminology, a stable log surface is a projective surface pair with semi-log-canonical singularities for which $K_X+\Delta$ is an ample $\Q$-Cartier divisor. The record does not specify the ground field, the allowed coefficients of the effective boundary $\Delta$, or additional conventions defining the class of surfaces. These choices must be fixed before a uniform bound can be stated unconditionally. On $U$, write $\omega_U$ for the invertible dualizing sheaf and interpret $rK_U$ as the line bundle $\omega_U^{\otimes r}$.

The first supplied question asks for the optimal positive integer $r$ such that $rK_U$ induces a birational map or an embedding. Under the interpretation that the bound is uniform, it asks for the least positive integers $r_{\mathrm{bir}}$ and $r_{\mathrm{emb}}$, separately, such that for every stable log surface in the intended class the map induced by $r_{\mathrm{bir}}K_U$ is birational onto its image, respectively the map induced by $r_{\mathrm{emb}}K_U$ is an embedding. Existence of such uniform integers is part of the question.

This formulation remains conditional: the record does not specify whether sections are taken on $U$ or obtained by restriction from $X$, whether the maps are required only on $U$ or on $X$, or whether an embedding means a locally closed immersion or a closed immersion. Since $U$ need not be proper, using all sections on $U$ does not automatically specify a finite-dimensional projective linear system. Moreover, the first question explicitly uses $K_U$, whereas the second uses $K_X+\Delta$; no identification or replacement of these divisors is justified by the supplied record.

For the second question, if $I$ denotes the global Cartier index of $K_X+\Delta$, define
\[
I=I(X,\Delta):=\min\{m\in\Z_{>0}:m(K_X+\Delta)\text{ is Cartier}\}.
\]
With this interpretation, the requested existence statement is
\[
\exists\,(X,\Delta)\text{ a stable log surface in the intended class}
\quad\text{such that}\quad
\mathcal O_X\bigl(5I(X,\Delta)(K_X+\Delta)\bigr)
\text{ is not very ample}.
\]
Here a line bundle is very ample if its complete linear system defines a closed immersion into projective space. The meaning of $I$ is not explicitly supplied, so this precise interpretation requires confirmation from the original source.
\subsection{Short English statement}
Determine the smallest uniform pluricanonical multiples giving birational maps and embeddings for the stated stable log surfaces, and ask whether the index-normalized log canonical multiple $5I(K_X+\Delta)$ can fail to be very ample. The ground field, boundary conventions, section spaces, embedding target, and meaning of $I$ require clarification. This is a statement-only draft requiring human mathematical review.
\subsection{Sources}
\begin{itemize}
\item Local registry record: permanent UnsolvedMath ID $30002343$, source problem code \texttt{catalogue entry}; project registry \texttt{data/unsolvedmath/problems.json}, with \texttt{catalog/problems.json} designated as the supplementary registry. The supplied record is titled ``Optimal Pluricanonical Bounds for Stable Log Surfaces.''
\item Supplied report citation: \emph{Complex Algebraic Geometry (2014)}, Oberwolfach Reports, DOI \texttt{10.4171/owr/2013/27}, pp.\ 1563--1627. \url{https://doi.org/10.4171/owr/2013/27}. The record gives proposed year $2013$ and names no individual proposer. The report-year discrepancy is unresolved.
\item Supplied literature reference: Liu--Rollenske, \emph{Pluricanonical maps of stable log surfaces}. \url{https://arxiv.org/abs/1211.1291}.
\item Supplied publication reference: \emph{Advances in Mathematics} publication. \url{https://doi.org/10.1016/j.aim.2014.03.009}.
\item Literature assessment dated exactly \texttt{2026-08-21}: the supplied curation labels the problem open and reports that the index multiples $4I$ and $8I$ are respectively base-point free and very ample, with stronger results for semi-canonical surfaces. It states that the optimal uniform $r$ and the proposed $5I$ counterexample question remain unresolved. This is supplied dated curation derived from the OWR record, not an independently verified literature assessment; its hypotheses, terminology, bounds, and time-sensitive status require independent checking.
\item The supplied statement status is \texttt{corrected\_verified}, accompanied by a claim of manual semantic, typographic, and LaTeX review. This records the supplied editorial assessment and does not establish independent source verification or mathematical correctness of this draft.
\end{itemize}
\end{document}
